\subsection{O que limita a resolução}
\label{sec:limites}

No sorteio nominal do teste, o RMSE em $z$ do baricentro calibrado cai de \num{55.18} para \qty{37.02}{\cm} entre as eficiências $\varepsilon$ de \num{0.005} e \num{1}, e de \num{42.45} para \qty{35.24}{\cm} entre as energias de \num{10} e \qty{100}{\MeV} (a Figura~\ref{fig:baricentro_luz} mostra médias). Com \qty{177.8}{\pe\per\MeV} em média, perto do rendimento médio do APEX, fica em \qty{38.23}{\cm} e, abaixo de $z = {}$\qty{20}{\m}, com menos luz, em \qty{13.79}{\cm}, próximo dos \qty{13.7}{\cm} do APEX a \qty{10}{\MeV} (Tabela~\ref{tab:referencias}). Detectar todos os fótons reduziria o erro do baricentro em \qty{5.43}{\cm} (\qty{4.96}{\cm} na média dos sorteios), e mais ainda sem o arredondamento da Equação~\eqref{eq:luz}, que a \qty{10}{\MeV} zera os canais mais fracos e eleva esse erro; por isso, com a mesma luz detectada, a varredura em energia fica abaixo da varredura em eficiência. No conjunto do teste, o erro do baricentro é sobretudo sistemático: com $\varepsilon = {}$\num{0.02}, o GATv2\discretionary{-}{-}{-}cascata (\qty{20.46}{\cm}) já erra menos em $z$ que o baricentro com todos os fótons. Essa parte sistemática vem dos vértices a partir de \qty{20}{\m}, onde o baricentro fica em \qty{63.16}{\cm} mesmo com todos os fótons (Seção~\ref{sec:onde_erro}); abaixo de \qty{20}{\m}, seu erro é sobretudo estatístico: de \qty{23.29}{\cm}, perto dos \qty{22.54}{\cm} do GATv2\discretionary{-}{-}{-}cascata, cai para \qty{8.33}{\cm} com todos os fótons.

O erro do GATv2-cascata também cai com mais luz (Seção~\ref{sec:onde_erro}), mas a luz diminui longe das paredes longas, e os modelos só foram treinados com \qty{10}{\MeV} e $\varepsilon = {}$\num{0.02}: para eles, nenhuma varredura separa o efeito da estatística de fótons do da posição. A influência da cobertura da grade não foi isolada. Com o treino completo, todo vértice de teste tem um de treino a no máximo \qty{50}{\cm}; com um quarto dos eventos, como na primeira comparação (Seção~\ref{sec:dificuldades}), \qty{27.6}{\percent} dos vértices de teste ficam sem esse vizinho, e a distância média sobe de \num{49.31} para \qty{56.18}{\cm}. Os modelos leem ainda \num{33} escalares do evento construídos à mão, e, na configuração final, nenhum experimento separou deles a contribuição da passagem de mensagens: o resultado é dos codificadores em grafos com esses escalares, não da passagem de mensagens nem da representação em grafos isoladamente.
